% ==================================================================
% id:       func-004
% topic:    פונקציות
% subtopic: תחום הגדרה ותכונות של פונקציות
% solution: full
% ==================================================================
הסבירו את ההבדל בין המושג
\underline{טווח של פונקציה}
ובין המושג
\underline{תמונה של פונקציה}.
תנו דוגמה.
\hint{חשבו על פונקציה שקבוצת היעד שלה גדולה יותר מאוסף הערכים שהיא מקבלת בפועל.}
\begin{solbox}
תהי
$f: A \rightarrow B$
פונקציה.
{\bf
הטווח
}
של $f$ הוא הקבוצה $B$.

{\bf
התמונה
}
של $f$ היא אוסף האיברים מתוך $B$ שהם תמונות של איברים ב-$A$.

דוגמה: נגדיר
$f: \rr \rightarrow \rr$
על ידי
$f(x) = x^2$.
אז הטווח של הפונקציה היא כל המספרים הממשיים. התמונה של הפונקציה היא אוסף המספרים הממשיים שהם ריבוע של מספר ממשי כלשהו, ולכן התמונה שווה ל-
$[0, \infty)$.
נשים לב שהתמונה תמיד מוכלת בטווח, אבל אינה בהכרח שווה לו. (מבחינה זאת הטווח הוא מושג מלאכותי במידה מסויימת)
\end{solbox}
